still integral over @@M0@@, with fibers of dimension @@M1@@; by 4.3, we may replace @@M2@@ by @@M3@@. Finally, by 4.3 and 3.10 (using resolution of singularities for surfaces), after replacing @@M4@@ by a flat covering of a nonempty open of @@M5@@, we may replace @@M6@@ by a smooth @@M7@@-scheme.

Now, by 4.6 (i) @@M8@@ (iii) and the finiteness and specialization theorem for \'etale cohomology groups (SGA 4 XVI 2.2), 5.1 follows from the next proposition.

\sgaStableTarget{sga6.e13.prop.5-2}\textbf{Proposition 5.2.} \emph{Let @@M9@@ be an algebraically closed field and @@M10@@ a proper smooth @@M11@@-scheme. Then the group @@M12@@ of algebraic cycles of codimension @@M13@@ on @@M14@@, modulo numerical equivalence, is free of rank @@M15@@, the Betti number.}

Indeed, if we forget the action of the roots of 1, \'etale cohomology @@M16@@ where @@M17@@ becomes a Weil cohomology: it satisfies Poincar\'e duality and the K\"unneth formula; and it receives a functorial, nonzero homomorphism from the ring of algebraic cycles modulo rational equivalence @@M18@@ (see X 7.13).

Let @@M19@@, where @@M20@@, be elements that map to a basis of the @@M21@@-vector subspace of @@M22@@ generated by algebraic cycles. Consider the homomorphism
@@M23@@
defined by @@M24@@, where @@M25@@ denotes the intersection number. Then one readily sees that the image of @@M26@@ is @@M27@@.

